%% ma: L12-B18-0004
%% lop: 12
%% chuong: 6
%% bai: 18
%% dang: 12.18.4
%% loai: DS
%% muc: [NB, NB, TH, TH]
%% dap_an: "ĐĐĐĐ"
%% nguon: "(NBV)-ĐỀ SỐ 8-MỖI NGÀY 1 ĐỀ THI 2025.pdf (D08, MỖI NGÀY 1 ĐỀ - Đề số 8), Phần II Câu 4"
%% kien_thuc: "Bài 18 (xác suất có điều kiện, công thức nhân xác suất)"
%% trang_thai: cho-duyet
%% ly_do: "Dạng 12.18.4 (xác suất có điều kiện và công thức nhân từ các tỉ lệ phần trăm, Đúng/Sai) là dạng mới đề xuất."
%% ghi_chu: "Đáp án gốc Đ Đ Đ Đ đúng. Ảnh gốc (hình minh họa dạ dày và stress) là ảnh trang trí, đã bỏ."
\begin{ex}
Khi kiểm tra sức khoẻ tổng quát của bệnh nhân ở một bệnh viện, người ta được kết quả như sau:
\begin{itemize}[leftmargin=1.5em,nosep]
  \item Có $40\%$ bệnh nhân bị đau dạ dày.
  \item Có $30\%$ bệnh nhân thường xuyên bị stress.
  \item Trong số các bệnh nhân thường xuyên bị stress có $80\%$ bệnh nhân bị đau dạ dày.
\end{itemize}
Chọn ngẫu nhiên $1$ bệnh nhân.
\choiceTF{\True Xác suất chọn được bệnh nhân thường xuyên bị stress là $0{,}3$}{\True Xác suất chọn được bệnh nhân bị đau dạ dày, biết bệnh nhân đó thường xuyên bị stress, là $0{,}8$}{\True Xác suất chọn được bệnh nhân vừa thường xuyên bị stress vừa bị đau dạ dày là $0{,}24$}{\True Xác suất chọn được bệnh nhân thường xuyên bị stress, biết bệnh nhân đó bị đau dạ dày, là $0{,}6$}
\loigiai{Gọi $A$: "bệnh nhân thường xuyên bị stress", $B$: "bệnh nhân bị đau dạ dày". Khi đó $P(A)=0{,}3$, $P(B)=0{,}4$, $P(B\mid A)=0{,}8$.

a) Đúng: $P(A)=0{,}3$. b) Đúng: $P(B\mid A)=0{,}8$.

c) Đúng: $P(AB)=P(A)\cdot P(B\mid A)=0{,}3\cdot0{,}8=0{,}24$.

d) Đúng: $P(A\mid B)=\dfrac{P(AB)}{P(B)}=\dfrac{0{,}24}{0{,}4}=0{,}6$.}
\end{ex}
